\subsection{RQ2: Metrics for assessing dataset quality}
\label{sec:rq2}
\draftnote{For each retained candidate add its definition, direction,
citation, implementation, assumptions, inputs and evidence of its
relationship to downstream utility. This table is an extraction framework,
not a recommendation to use every measure.}

\begin{table}[htbp]
\caption{Candidate dimensions and measures to investigate.}
\label{tab:dimensions}
\small
\renewcommand{\arraystretch}{1.2}
\begin{tabularx}{\linewidth}{@{}p{25mm}YY@{}}
\toprule
\textbf{Dimension} & \textbf{Candidate measures} & \textbf{Interpretation questions} \\
\midrule
Realism & Human ratings; physical and semantic validity checks. & Are criteria explicit and raters consistent? \\
Similarity & Fr\'echet Inception Distance (FID); Kernel Inception Distance (KID); class-conditional comparisons. & Which features, reference data and preprocessing? \\
Coverage & Generative precision/recall; $\alpha$-Precision and $\beta$-Recall; scenario occupancy. & Are rare cases and joint conditions represented? \\
Bias & Representation ratios; subgroup error gaps; worst-group performance. & Which population and group definitions? \\
Labels & Audited label error; box/mask agreement. & Are reference annotations independently checked? \\
Utility & Classification accuracy or macro-F1; detection mAP; segmentation mIoU. & Performance on independent real test data? \\
Authenticity & Authenticity; duplicate and nearest-neighbour audits. & Is generator training data available? \\
\bottomrule
\end{tabularx}
\end{table}

\subsubsection{Validity and limitations}
\draftnote{Examine sample-size effects, feature dependence, uncertainty
and implementation comparability. Distinguish generative precision/recall
from classifier precision/recall. Investigate whether paired measures such
as SSIM or PSNR fit the data; avoid arbitrary real/synthetic pairing.
Distinguish authenticity diagnostics from formal privacy guarantees.}

\subsubsection{Selecting complementary metrics}
\draftnote{Use evidence to propose a task-dependent set of distributional,
coverage, label, subgroup and utility measures. Explain what each adds
and misses, and whether acceptance thresholds are supported.
End with a direct answer to RQ2.}
